%! Unit 1, Lessons 1.0–1.4 — ACTUAL Quiz (Answer Key)
% The scored mid-unit quiz. Item for item it is the parallel form of ../sample_quiz/: same
% sections, same six multiple-choice targets in the same order, same two free-response sets with
% the same parts — every context different, so a student who worked the sample has practised the
% moves and not memorised the answers. The quantitative display is deliberately a HISTOGRAM here
% where the sample used a dotplot: Lesson 1.3 asks students to read either.
% NAMESTRIP: \namedateperiod mirrors the blank — structural, so the pages stay in lockstep.
% NO teachernote here — scoring guidance belongs in the teacher's copy, never in a key.
% GENERATED FROM the blank. PAGE LOCKSTEP with ../../quiz_keys/actual_quiz_key/: every \writelines{n} here is answered by
% exactly n \ansline{} there, each terse enough not to wrap. Nothing else differs between the
% two files except the preamble, the header, and the \ans{} marks on the correct MC options.
\documentclass[12pt]{article}
\usepackage{apstats-article}
\usepackage{apstats-key}   % pulls in -boxes; do NOT also load -boxes

\begin{document}

\pageheader{Unit 1 \quad Lessons 1.0--1.4}{Quiz --- Answer Key}

\namedateperiod
\vspace{0.10in}

\boxguard[24]
\begin{scenariobox}[The Harbor Aquarium]{navy}
\small
The Harbor Aquarium keeps one row of records for each of the \textbf{60 guided tours} it ran last
month. Four of the 60 rows are shown below.

\vspace{6pt}
\begin{center}
\renewcommand{\arraystretch}{1.2}
\begin{tabular}{@{} l l c r @{}}
\rowcolor{sky}
\textbf{Tour} & \textbf{Theme} & \textbf{Tour code} & \textbf{Visitors} \\
\midrule
Tuesday 10:00  & Tide pools  & T-07 & 24 \\
Tuesday 13:30  & Sharks      & T-19 & 31 \\
Thursday 09:15 & Penguins    & T-28 & 18 \\
Saturday 11:00 & Coral reef  & T-46 & 27 \\
\end{tabular}
\end{center}

\vspace{6pt}
Across all 60 tours, the themes are: \textbf{Tide pools 21}, \textbf{Sharks 15},
\textbf{Penguins 15}, \textbf{Coral reef 9}.
\end{scenariobox}

\vspace{0.14in}

\boxguard[16]
\begin{headlinebox}{goldbg}
\small
\textbf{Section I --- Multiple Choice.} Circle the single best answer. Read all five options ---
more than one of them is about the context, not just the arithmetic.
\end{headlinebox}

\vspace{0.10in}

\begin{enumerate}[leftmargin=1.9em, itemsep=0.13in]

  \item In the aquarium's records, the \textbf{individuals} are
        \begin{enumerate}[label=\textbf{(\Alph*)}, leftmargin=2.2em, itemsep=2pt]
          \item the four tour themes.
          \item \ans{the 60 guided tours.}
          \item the visitors who took the tours.
          \item the aquarium itself.
          \item the numbers of visitors.
        \end{enumerate}

  \item Which of the following is a \textbf{statistical question} about these tours?
        \begin{enumerate}[label=\textbf{(\Alph*)}, leftmargin=2.2em, itemsep=2pt]
          \item How many visitors took tour T-19?
          \item What is the theme of tour T-28?
          \item How many tours did the aquarium run last month?
          \item \ans{How much does the number of visitors vary across the 60 tours?}
          \item Does the aquarium have penguins?
        \end{enumerate}

  \item \textbf{Tour code} is best classified as
        \begin{enumerate}[label=\textbf{(\Alph*)}, leftmargin=2.2em, itemsep=2pt]
          \item quantitative, because part of each value is a number.
          \item quantitative and discrete, because the codes run from 01 to 60.
          \item quantitative and continuous, because codes could be subdivided.
          \item \ans{categorical, because the code labels a tour rather than recording an amount.}
          \item categorical, because it begins with a letter.
        \end{enumerate}

  \item Bayview Aquarium ran 30 tours last month, 12 of them tide-pool tours. Which statement
        about the two aquariums is correct?
        \begin{enumerate}[label=\textbf{(\Alph*)}, leftmargin=2.2em, itemsep=2pt]
          \item Harbor is more tide-pool-focused, because 21 is greater than 12.
          \item \ans{Bayview is more tide-pool-focused, because 40\% is greater than 35\%.}
          \item The two aquariums are equally tide-pool-focused.
          \item No comparison is possible, because the aquariums ran different numbers of tours.
          \item Harbor is more tide-pool-focused, because it ran more tours overall.
        \end{enumerate}

  \item \textbf{Number of visitors on a tour} is
        \begin{enumerate}[label=\textbf{(\Alph*)}, leftmargin=2.2em, itemsep=2pt]
          \item categorical, because tours are grouped by size.
          \item continuous, because a tour could in principle be any size.
          \item \ans{discrete, because it is a count and 24.5 visitors is impossible.}
          \item continuous, because the values are whole numbers.
          \item neither, because a count cannot be graphed.
        \end{enumerate}

  \item A histogram of tour lengths piles up at the low end and trails off toward the high end,
        with a few tours far to the right. This distribution is
        \begin{enumerate}[label=\textbf{(\Alph*)}, leftmargin=2.2em, itemsep=2pt]
          \item skewed left, because most of the tours sit on the left.
          \item \ans{skewed right, because the longer tail runs out to the right.}
          \item symmetric, because it has one peak.
          \item uniform, because every tour has some length.
          \item impossible to name without the center.
        \end{enumerate}

\end{enumerate}

\vspace{0.14in}

% No \newpage here: the guard keeps the Section II strip off a page edge, and letting
% the section flow saves the quiz a near-empty final page.
\boxguard[10]
\begin{headlinebox}{goldbg}
\small
\textbf{Section II --- Free Response.} Show your reasoning. Every answer must be written
\emph{in the context of the problem} --- name the variable and its units.
\end{headlinebox}

\vspace{0.10in}

\begin{enumerate}[leftmargin=1.9em, itemsep=0.16in, resume]

  \item Use the Harbor Aquarium records above.

        \vspace{6pt}
        \textbf{(a)} Identify the individuals, then name one categorical variable and one
        quantitative variable in the table.\\[4pt]
        \ansline{Individuals: the 60 guided tours, one per row.}\\
        \ansline{Categorical: theme. Quantitative: number of visitors on the tour.}

        \vspace{6pt}
        \textbf{(b)} Give the relative frequency of each of the four themes, to the nearest tenth
        of a percent.\\[4pt]
        \ansline{Tide pools 35.0\%, Sharks 25.0\%, Penguins 25.0\%, Coral reef 15.0\%.}

        \vspace{6pt}
        \textbf{(c)} A newsletter writes, ``Harbor is the region's tide-pool aquarium --- it ran
        more tide-pool tours than Bayview did.'' Harbor ran 21 of 60; Bayview ran 12 of 30.
        Explain what is wrong with that reasoning, and state which aquarium a share comparison
        favours.\\[4pt]
        \ansline{Different sizes, so the counts cannot be compared: 60 tours give Harbor}\\
        \ansline{more of everything. Shares are the fair number: 21/60 = 35.0\% against}\\
        \ansline{12/30 = 40.0\%, so a share comparison favours Bayview.}

        \vspace{6pt}
        \textbf{(d)} Write one statistical question that could be answered using these records,
        and say which of your two checks makes it statistical.\\[4pt]
        \ansline{``How much does the number of visitors vary across the 60 tours?''}\\
        \ansline{It anticipates variability and is answered by collecting data.}

  \item The histogram below shows the \textbf{length of visit}, in minutes, of all 40 visitors
        recorded on one Saturday.

        \vspace{4pt}
        \noindent\hfil
        \begin{tikzpicture}[scale=0.75]
          \foreach \i/\c in {0/3, 1/7, 2/12, 3/9, 4/5, 5/2, 6/0, 7/2} {
            \fill[navy] ({\i*0.9},0) rectangle ({(\i+1)*0.9},{\c*0.16});
            \draw[white, very thin] ({\i*0.9},0) rectangle ({(\i+1)*0.9},{\c*0.16});}
          \draw[->, thick] (0,0) -- (0,2.4);
          \draw[->, thick] (0,0) -- (7.8,0);
          \foreach \y/\lab in {0/0, 0.64/4, 1.28/8, 1.92/12} {
            \draw (-0.07,\y) -- (0.07,\y); \node[left, font=\tiny, xshift=-1pt] at (-0.07,\y) {\lab};}
          \foreach \i/\lab in {0/20, 1/30, 2/40, 3/50, 4/60, 5/70, 6/80, 7/90, 8/100} {
            \draw ({\i*0.9},-0.07) -- ({\i*0.9},0.07);
            \node[below, font=\tiny] at ({\i*0.9},-0.1) {\lab};}
          \node[font=\tiny] at (3.6,-0.8) {Length of visit (minutes)};
          \node[font=\tiny, rotate=90] at (-0.62,1.2) {Number of visitors};
        \end{tikzpicture}\hfil
        \vspace{6pt}

        \textbf{(a)} Is length of visit discrete or continuous? Justify from the variable itself,
        not from the display.\\[4pt]
        \ansline{Continuous --- length of visit is a measured amount of time.}\\
        \ansline{Between 42 and 43 minutes, 42.4 minutes is possible; rounding changes nothing.}

        \vspace{6pt}
        \textbf{(b)} Describe the distribution of length of visit completely.\\[4pt]
        \ansline{Skewed right: the bars pile up at 40--50 minutes, tail out to the right.}\\
        \ansline{A gap at 80--90 minutes, and a cluster of two visitors at 90--100.}\\
        \ansline{Center is in the 40--50 minute block; visits run from about 20 to 100}\\
        \ansline{minutes --- a great deal of variability among the 40 visitors.}

        \vspace{6pt}
        \textbf{(c)} A classmate writes, ``It is skewed left, because most of the bars are on the
        left side.'' Correct the classmate, and say what actually decides the name.\\[4pt]
        \ansline{A skew is named for the side the longer tail runs out on, not the}\\
        \ansline{side the data piles up on. This tail reaches long, so: skewed right.}

        \vspace{6pt}
        \textbf{(d)} The aquarium's director proposes deleting the two visitors in the 90--100
        minute block before reporting a typical visit, calling them ``clearly errors.'' Give one
        reason that is a poor decision, and say what should be done instead.\\[4pt]
        \ansline{Nothing marks them as errors --- a visitor really can stay 95 minutes.}\\
        \ansline{Investigate those two visits and report them as part of the description.}

\end{enumerate}

\end{document}
